\track{T6\quad Langlands 对应、$\mathcal D$-模与几何背景}{本轨道 21 条：终点站——物理侧的 A/B 膜范畴被翻译成 $\mathcal D$-模，几何 Langlands 对应就是这两个范畴的等价。}

\begin{word}{几何 Langlands}{Geometric Langlands}{geometric-langlands}{§1 p.2}
\mathdef{设 $C$ 为紧黎曼面，$^LG$ 为 Langlands 对偶群。几何 Langlands 对应断言：$C$ 上平坦 $^LG$-丛的模空间 $\mathrm{Loc}_{^LG}(C)$ 上相干层的导出范畴，与 $\mathrm{Bun}_G(C)$ 上扭化 ${\cal D}$-模的导出范畴等价（准确到一个 $Z_2$-gerbe 的扭化），把点 $E$（一个局部系统）送到以 $E$ 为特征值的 Hecke 特征层 $F_E$；一般 $G$ 的「一一对应」须换成 Hecke 特征性 $H_V(F_E)\cong V_E\boxtimes F_E$。}
\mathex{$G=GL_1$ 时 $\mathrm{Bun}_{GL_1}=\mathrm{Pic}(C)$，对应退化为 Fourier--Mukai 变换，是类域论的几何版本；$C=\mathbb P^1$ 时 $\pi_1$ 平凡，$\mathrm{Loc}$ 只有平凡局部系统，Hecke 特征层相应退化；$E$ 为椭圆曲线、$G=SL_2$ 时，平坦 $SL_2$-丛由 $\pi_1(E)=\mathbb Z^2$ 的两个交换矩阵给出，$\mathrm{Bun}_{SL_2}$ 连通（$\pi_1(G)=0$）。}
\physdef{物理侧：把四维 $N=4$ 超对称规范理论沿 $\Sigma\times C$ 紧化并作拓扑扭化，$G$ 侧得到靶空间为 Hitchin 模空间 $\mathcal M_H(G,C)$ 的 $A$-模型，$^LG$ 侧得到靶空间为 $\mathcal M_H(^LG,C)$ 的 $B$-模型；S-对偶交换 $G$ 与 $^LG$，从而交换 $A$-膜与 $B$-膜。}
\physex{电特征膜是 $B$-模型里支撑在点 $E$ 的 $0$-膜；其 S-对偶是缠在 Hitchin 纤维上、带平坦 Chan-Paton 丛的 $A$-膜，即磁特征膜。纤维上的 T-对偶给出这对膜的全部内容，也给出两个 Hitchin 纤维化底空间的自然等同。}
\bridge{两边的对象：$A$-膜 $\leftrightarrow$ $\mathrm{Bun}_G(C)$ 上的扭化 ${\cal D}$-模；$B$-膜（相干层）$\leftrightarrow$ $\mathrm{Loc}_{^LG}(C)$ 上的对象。特征层 $\leftrightarrow$ 特征膜：$E$ 给出 $B$ 侧的 structure sheaf $\mathcal O_E$（电特征膜），其镜像 $A$-膜缠在 $\mathcal M_H(G,C)$ 的一条 Hitchin 纤维上，再经「$A$-膜 $\leftrightarrow$ ${\cal D}$-模」译成 $F_E$；Hecke 变换（线算子）作用在 $F_E$ 上并使其成为特征层。自由参数是 $E$ 的共轭类；原文见 §1 p.2 与 §11。}
\intuit{像傅里叶变换：一侧是点（局部系统），另一侧是波（特征层），Hecke 算子则是频率平移。边界：这是非分歧、非奇异情形的陈述，且范畴等价比「每个对象配一个对象」更强，不能逐条读成集合对应。}
\pitfall{把几何 Langlands 当成数论 Langlands 的自守形式对应直接照搬：几何侧的对象是层与范畴，不是函数，没有 L-函数与 Euler 积；照搬「每个局部系统配一个自守表示」会推出不存在的一一对应，并漏掉 packet 现象。}
\end{word}

\begin{word}{Langlands 纲领}{Langlands Program}{langlands-program}{§1 p.2}
\mathdef{数域 $F$ 上，Langlands 纲领预测 $n$ 维 Galois 表示与 $\mathrm{GL}_n$ 的自守表示之间的对应（无分歧处对应，分歧处带局部参数）；一般约化群 $G$ 要用 L-包（L-packet，同一参数下的一族表示）与函子性来表述。核心结构是 Langlands 对偶群 $^LG$，它交换 weight 格与 coweight 格、单根与单余根。}
\mathex{$n=1$ 即类域论（Artin 互反律）；$F=\mathbb Q,n=2$ 对应模形式；函数域 $\mathbb F_q(X)$ 情形由 Drinfeld（$n=2$）与 L. Lafforgue（一般 $n$）证明；$G=\mathrm{SL}_2$ 用 $^LG=\mathrm{PGL}_2$ 的参数化，$E_8$ 自对偶。一般群由 Arthur 的 $\mathrm{SL}_2$ 得到 L-包，故不是一一对应。}
\physdef{物理侧：GNO 对偶指出规范群 $G$ 与其对偶 $^LG$ 的电、磁通量格互换；在 $N=4$ 超对称理论中 Montonen--Olive 对偶连同 theta 角生成 $SL(2,\mathbb Z)$，其「双群」正是 $^LG$（Atiyah 1977 年的提示，见原文 §1）。}
\physex{$\mathrm{SU}(2)$ 与 $\mathrm{SO}(3)$、$\mathrm{Sp}(2n)$ 与 $\mathrm{SO}(2n+1)$ 是标准例子；'t Hooft 线算子按 $^LG$ 的表示分类，Wilson 线按 $G$ 的表示分类，二者在 S-对偶下互换，电荷与磁荷的角色也对调。}
\bridge{对应：GNO 把电通量格 $X_*(T)$ 与磁通量格 $X^*(T)$ 互换，正是 Langlands 对偶交换 weight/coweight 格；数论侧的「Galois 表示 $\leftrightarrow$ 自守表示」在几何侧被换成「局部系统 $\leftrightarrow$ Hecke 特征层」，自守表示的几何替代是 $\mathrm{Bun}_G$ 上的自守层。原文 §1 p.2 只借用对偶群结构。}
\intuit{像同一份物理的两种变量：电的语言与磁的语言互为傅里叶对偶；数学上则是「群」与「其表示论」的对偶。边界：物理上 S-对偶只在特定理论中成立，数学上一般群仍是猜想，L-包说明对应不是一一的。}
\pitfall{把 Langlands 对偶当成群同构：$\mathrm{SL}_2$ 的对偶是 $\mathrm{PGL}_2$，并不相等；把一般群的对应当成一一对应（忽略 L-包与函子性）；或把几何 Langlands 的层对应误当作数论纲领的逐字翻译。}
\end{word}

\begin{word}{D-模}{D-module}{d-module}{§11 p.184}
\mathdef{复流形 $X$ 上，微分算子层 ${\cal D}_X$ 由 ${\cal O}_X$ 与切向量场生成，满足 Leibniz 规则；左 ${\cal D}_X$-模即带联络的 ${\cal O}_X$-模 $M$：$\nabla:M\to M\otimes\Omega^1_X$，$\nabla(fs)=f\nabla s+s\otimes df$，并满足可积性 $\nabla^2=0$。扭化版本见「扭化 D-模」条。}
\mathex{平坦向量丛 $(E,\nabla)$ 的截面层是 ${\cal D}_X$-模；$X=C$ 时 ${\cal D}_C$ 局部由 ${\cal O}_C$ 与 $\partial_z$ 生成；$G=GL_1$ 时 $\mathrm{Bun}_{GL_1}=\mathrm{Pic}(C)$，其上的 ${\cal D}$-模由 Fourier--Mukai 变换分类。$\delta$-函数模、局部上同调模等非 ${\cal O}_X$-相干对象也是 ${\cal D}$-模。}
\physdef{物理侧：${\cal D}_X$-模是余切丛 $T^*X$ 上 $B$-膜（相干层）的非交换形变；开弦两端落在同一张膜上时其代数就是 ${\cal D}_X$，膜上的态就是 ${\cal D}_X$-模。这就是 Kapustin 的「$A$-膜 $\leftrightarrow$ ${\cal D}$-模」翻译器的代数基础。}
\physex{$T^*C$ 零截面上的 $B$-膜给出 $C$ 上的 ${\cal D}$-模；典范余各向同性膜（c.c. 膜）的开弦代数正是 ${\cal D}_{K^{1/2}}$，故几何 Langlands 里出现的是扭化 ${\cal D}$-模而非普通 ${\cal D}$-模。}
\bridge{对应：符号分次 $\mathrm{gr}{\cal D}_X\cong\mathrm{Sym}\,T_X$ 给出余切丛上的函数代数，故 ${\cal D}_X$-模是「量子化余切丛」上的对象；${\cal D}$-模的特征簇（characteristic variety）是其 $\mathrm{gr}$ 的支撑，与 $A$-膜在 $\mathcal M_H$ 中的支撑相呼应。原文 §11 p.184。}
\intuit{${\cal D}$-模把「微分方程」变成几何对象：一个 ${\cal D}$-模是一组解空间的代数封装，像量子力学中 $[x,p]=i\hbar$ 的表示。边界：它不是希尔伯特空间，没有内积与态。}
\pitfall{把 ${\cal D}$-模等同于「带平坦联络的向量丛」：后者只是相干且 ${\cal O}$-相干的特例，而 ${\cal D}$-模可以是无限秩、可以是 $\delta$-函数型；也把 ${\cal D}$-模当成「量子化之后的代数」，原文指的是微分算子层的模。}
\end{word}

\begin{word}{微分算子层}{Sheaf of Differential Operators}{sheaf-of-differential-operators}{§1 p.5；§11.2 p.196}
\mathdef{$X$ 上由 ${\cal O}_X$ 与一阶微分（向量场）生成的环层 ${\cal D}_X$：局部坐标下 ${\cal D}_X={\cal O}_X[\partial_1,\dots,\partial_n]$，$[\partial_i,f]=\partial f/\partial x^i$。它是非交换的，带滤过，$\mathrm{gr}{\cal D}_X\cong\mathrm{Sym}\,T_X$。作用在线丛 $L$ 的截面上的版本记 ${\cal D}_L$，且 ${\cal D}_L^{\mathrm{op}}\cong{\cal D}_{L^{-1}\otimes K_X}$。}
\mathex{$X=C$ 时 ${\cal D}_C$ 是 $C$ 上全纯微分算子层；$X=\mathbb P^1$ 的整体一阶算子给出 $\mathfrak{sl}_2$ 作用；取 $L=K_C^{1/2}$ 得 ${\cal D}_{K^{1/2}}$，这正是几何 Langlands 所需的 c.c. 弦代数（原文 §11.2 p.196）。}
\physdef{物理侧：${\cal D}_X$ 是端点都落在同一张膜上的开弦代数；$\mathrm{gr}{\cal D}_X$ 是 $T^*X$ 上的函数，对应经典极限与算子符号。扭化版本来自让开弦端点带一个线丛（Chan-Paton 线丛）。}
\physex{c.c. 膜的开弦代数必是 ${\cal D}_{K_M^{1/2}}$：时间反演要求代数同构于其相反代数，而 ${\cal D}_L^{\mathrm{op}}\cong{\cal D}_{L^{-1}\otimes K_M}$，故 $L=K_M^{1/2}$；$M$ 自旋时 $K_M^{1/2}$ 是普通线丛。}
\bridge{对应：符号映射 $\sigma:{\cal D}_X\to\mathrm{Sym}\,T_X$ 就是经典--量子对应；${\cal D}_X$-模的特征簇是 $\sigma$ 的支撑，物理上即 $B$-膜在 $T^*X$ 中的支撑。原文 §1 p.5 与 §11.2 p.196。}
\intuit{${\cal D}_X$ 是「位置与动量的非交换代数」的层化版本，像量子力学里 $[x,p]=i\hbar$ 的代数逐点粘起来。边界：它只是余切丛上的形变，不是把底空间本身非交换化。}
\pitfall{把 ${\cal D}_X$ 当成交换环或「所有微分算子」的平凡层；把 ${\cal D}_L$ 与 ${\cal D}_X$ 混为一谈（差一个线丛，且相反代数对应 $L^{-1}\otimes K_X$，这正是 c.c. 膜要选 $K^{1/2}$ 的原因）。}
\end{word}

\begin{word}{扭化 D-模}{Twisted D-module}{twisted-d-module}{§11.2 p.196}
\mathdef{给定复流形 $X$ 与线丛 $L$，令 ${\cal D}_L$ 为作用在 $L$ 的局部截面上的一阶微分算子层（${\cal D}_L={\cal D}_X\otimes_{\cal O_X}L$）；${\cal D}_L$-模称 $L$-扭化 ${\cal D}$-模。更一般地可用扭化参数 $\lambda$ 记作 ${\cal D}_X^\lambda$；几何 Langlands 取 $L=K_X^{1/2}$，即原文的 ${\cal D}_{K^{1/2}}$。局部上 ${\cal D}_L$ 与 ${\cal D}_X$ 同构，整体上二者相差一个线丛（一个扭化类），故整体截面与 Hecke 作用都不同。}
\mathex{$X=C$ 上取 $L=K_C^{1/2}$ 得 ${\cal D}_{K_C^{1/2}}$；$X=\mathcal M$ 为 Hitchin 模空间时，c.c. 弦代数 ${\cal D}_{K_{\mathcal M}^{1/2}}$ 就是扭化 ${\cal D}$-模层的原型；$\mathrm{Bun}_G$ 上几何 Langlands 的特征层取值于临界扭化，故是扭化 ${\cal D}$-模。}
\physdef{物理侧：扭化 ${\cal D}$-模对应带背景线丛（或 't Hooft 通量）的开弦代数；普通 ${\cal D}$-模是「无背景通量」的特例。扭化参数对应背景 B-场/线丛的选取。}
\physex{$M$ 非自旋时 $K_M^{1/2}$ 不存在为普通线丛，须用 $Z_2$-gerbe，出现真正的扭化；$M$ 自旋时扭化可平凡化，但几何 Langlands 的层仍按 $K^{1/2}$ 的扭化书写。}
\bridge{对应：扭化类在物理侧是背景通量，在数学侧是线丛的转移数据；几何 Langlands 的 Hecke 特征层是 $\mathrm{Bun}_G$ 上按 $K^{1/2}$ 扭化的 ${\cal D}$-模，故「$A$-膜 $\to$ 扭化 ${\cal D}$-模」而非普通 ${\cal D}$-模。原文 §11.2 p.196。}
\intuit{像在磁场中做量子力学：普通导数要换成协变导数，算子代数整体上「拧」了一下。边界：局部看与普通 ${\cal D}$-模无异，差别只在整体截面与 Hecke 算子的作用。}
\pitfall{把扭化 ${\cal D}$-模当成普通 ${\cal D}$-模：局部同构会诱使人以为整体也一样，但 twisting 改变整体截面与 Hecke 特征性的表述；也忽略扭化类必须整体一致、不能逐点任意选取。}
\end{word}

\begin{word}{平坦联络}{Flat Connection}{flat-connection}{§3.2 p.25}
\mathdef{光滑或复流形 $X$ 上主 $G$-丛 $E$ 的联络 $\nabla:\Gamma(E)\to\Gamma(E\otimes\Omega^1)$，满足 $\nabla(fs)=f\nabla s+s\otimes df$；平坦即曲率 $F_\nabla=\nabla^2=0$。等价地，其局部常值截面层是局部系统；在复曲线 $C$ 上，全纯平坦 $G_{\mathbb C}$-丛的同构类由 $\mathrm{Hom}(\pi_1(C),G_{\mathbb C})/G_{\mathbb C}$ 给出（Riemann--Hilbert 对应）。}
\mathex{$C$ 上平坦 $SL_2$-丛由 $\pi_1(C)$ 的生成元 $A_i,B_i$ 给出，满足 $\prod_i[A_i,B_i]=1$；$G=GL_1$ 时 $\mathrm{Hom}(\pi_1(C),\mathbb C^*)$ 是 $2g$ 维环面 $(\mathbb C^*)^{2g}$；亏格 $g\ge 2$ 时表示簇的维数为 $(2g-2)\dim G$。}
\physdef{物理侧：平坦联络是场强为零的规范场 $A$；绕闭曲线的 holonomy 给出 Wilson 圈。$N=4$ 理论的 BPS 方程是 Hitchin 方程 $F-\varphi\wedge\varphi=0$、$D\varphi=D\star\varphi=0$，取 $\varphi=0$ 时退化为 $F=0$。}
\physex{$C$ 上的平坦 $U(1)$ 联络由 $H^1(C,\mathbb R)/H^1(C,\mathbb Z)\cong(S^1)^{2g}$ 参数化，绕洞的 Aharonov--Bohm 相位就是 holonomy；复平坦联络 $A+i\varphi$ 给出 $\mathcal M_H$ 的复结构 $J$。}
\bridge{对应：平坦联络 $=$ 局部系统 $=$ $\pi_1(C)$ 的表示；在 $\mathcal M_H$ 的复结构 $J$ 下，其点就是半稳定平坦丛，即几何 Langlands 左端的 $\mathrm{Loc}$。oper 是 de Rham 版本 $\mathcal Y(^LG,C)$ 中带 Borel 结构的半维子簇（一类特定全纯对象），不是一般平坦联络。原文 §3.2 p.25。}
\intuit{平坦联络是「局部没有曲率、整体有挠」的对象，像把纸条粘成莫比乌斯带：每小段都平，绕一圈却转了向。边界：平坦不等于平凡，$F=0$ 只保证局部可积。}
\pitfall{以为 $F=0$ 的联络都能整体写成 $g^{-1}dg$：这要求丛平凡或空间单连通；也把 oper 与平坦联络混同——oper 带有额外的 Borel/抛物结构，是 $\mathrm{Loc}$ 中的特殊子簇。}
\end{word}

\begin{word}{局部系统}{Local System}{local-system}{§10.4 p.171（原文用 locally constant）}
\mathdef{局部系统是局部常值层：每点附近同构于常值层 $\underline V$，等价于 $\pi_1(C,\mathrm{pt})$ 到 $\mathrm{GL}(V)$ 的表示，也等价于向量丛加平坦联络。其 Betti 版本的模空间 $\mathrm{Loc}_{^LG}(C)=\mathrm{Hom}(\pi_1(C),{}^LG)/{}^LG$ 是特征簇（character variety）；de Rham 版本则是平坦联络模空间。}
\mathex{$C=\mathbb P^1$ 时 $\pi_1$ 平凡，$\mathrm{Loc}$ 只有一个点（平凡局部系统）；$C=E$ 为椭圆曲线时 $\pi_1=\mathbb Z^2$，$G=GL_1$ 给出 $\mathrm{Hom}(\mathbb Z^2,\mathbb C^*)\cong\mathbb C^*\times\mathbb C^*$，即对偶环面；$G=SL_2$ 时同一底空间上的 $\mathrm{Loc}$ 是带奇点的二维复代数簇。}
\physdef{物理侧：局部系统是「电」的对象，其 Wilson 圈期望值由 monodromy 给出；在几何 Langlands 中它标号电特征膜——$B$-模型靶空间里支撑在点 $E$ 的 $0$-膜。}
\physex{点 $E$ 的对偶磁特征膜缠在 $\mathcal M_H(G)$ 的一条 Hitchin 纤维上，并带平坦 Chan-Paton 丛；纤维方向上的 T-对偶把 $\mathrm{Loc}_{^LG}$ 的点与它配对。}
\bridge{对应：$\mathrm{Loc}_{^LG}$ 是特征簇的 Betti 版本，平坦联络模空间 $\mathcal Y(^LG,C)$ 是 de Rham 版本；两者经 Riemann--Hilbert 对应，而 $\mathcal Y$ 又与 Higgs 丛模空间经 non-abelian Hodge 对应。原文 §10.4 p.171 用 locally constant，未用 local system 一词。}
\intuit{局部系统像「给每条闭曲线配一个矩阵，且只依赖同伦类」：绕洞的走法决定变换。边界：$\mathrm{Loc}$ 是栈而非流形，点带自同构，不能只数同构类。}
\pitfall{把局部系统当成普通向量丛（丢失平坦结构）；或把 $\mathrm{Loc}$ 与 $\mathcal M_H$ 混同：前者是 Betti 侧的表示簇，后者是超凯勒流形，两者只在复结构 $J$ 下并加半稳定性后才重合。}
\end{word}

\begin{word}{基本群}{Fundamental Group}{fundamental-group}{§5.6 p.77}
\mathdef{$\pi_1(C,\mathrm{pt})$ 是以闭曲线同伦类为元素、以连接为乘法的群；亏格 $g$ 的紧黎曼面有表现 $\langle a_1,b_1,\dots,a_g,b_g\mid\prod_i[a_i,b_i]=1\rangle$。它与局部系统的关系是 $\mathrm{Hom}(\pi_1(C),G)/G$；其交换化是 $H_1(C,\mathbb Z)\cong\mathbb Z^{2g}$。}
\mathex{$g=0$：$\pi_1$ 平凡，故 $C=\mathbb P^1$ 上只有平凡局部系统；$g=1$：$\pi_1=\mathbb Z^2$，$\mathrm{Hom}(\mathbb Z^2,\mathbb C^*)=(\mathbb C^*)^2$；$g\ge 2$：$\pi_1$ 非交换，$\mathrm{Hom}(\pi_1,SL_2)$ 模共轭后的维数是 $6g-6$。}
\physdef{物理侧：$\pi_1(C)$ 的同伦类正是 Wilson 圈可以绕的路线；绕洞一圈的 holonomy 是表示矩阵，'t Hooft 线算子则探测磁通量。S-对偶把 $\pi_1(C)$ 的表示与 $^LG$ 的表示配对。}
\physex{环面上两类基本圈 $a,b$ 生成 $\pi_1$，沿 $a,b$ 的 holonomy 矩阵交换即平坦性；Aharonov--Bohm 相位是 $\pi_1$ 表示的一维情形。}
\bridge{对应：$\pi_1(C)$ 到 $^LG$ 的同态 $\leftrightarrow$ $C$ 上平坦 $^LG$-丛 $\leftrightarrow$ 几何 Langlands 左端的电特征膜；GNO 对偶把 $\pi_1$ 的表示与 $^LG$ 的表示配对。原文 §5.6 p.77。}
\intuit{$\pi_1$ 记录「绕洞的走法」，表示记录「绕一圈后的变换」，像给每个洞装一扇旋转门。边界：$\pi_1$ 非交换，亏格大于零时表示不是自由参数，且须模共轭。}
\pitfall{用 $H_1$（交换化）代替 $\pi_1$：阿贝尔情形无妨，非阿贝尔时 $\mathrm{Hom}(\pi_1,G)$ 比 $\mathrm{Hom}(H_1,G)$ 小得多；也把同态本身当成点而忘记模共轭，得到维数虚高的「模空间」。}
\end{word}

\begin{word}{不可约表示}{Irreducible Representation}{irreducible-representation}{§5.4 p.73}
\mathdef{群或结合代数的表示 $V$ 若没有非平凡子表示，则称不可约。对 $\pi_1(C)$ 的表示，不可约等价于对应平坦丛没有 $\nabla$-不变子丛，也等价于 $\mathrm{Loc}$ 中的点不落在任何真抛物子群的轨道上。$G=GL_1$ 时所有表示都是一维，不可约性退化。}
\mathex{$G=SL_2$ 的平坦丛可约当且仅当存在不变线子丛，即 $A_i,B_i$ 可同时三角化（如对角常值情形）；$C=\mathbb P^1$ 上 $\pi_1$ 平凡，只有平凡表示；$E$ 为椭圆曲线时不可约平坦 $SL_2$-丛对应的特征层唯一。}
\physdef{物理侧：不可约性意味着规范场没有未破缺的子结构，膜没有连续稳定子；几何 Langlands 中 $E$ 不可约时 Hecke 特征层 $F_E$ 唯一且不可约、holonomic、具正则奇点。}
\physex{可约局部系统（如对角的平坦 $SL_2$-丛）在 Hitchin 纤维上不是孤立点，其「对偶特征膜」带连续参数；不可约点则给出孤立膜及其唯一对偶。}
\bridge{对应：$E$ 不可约 $\leftrightarrow$ $F_E$ 唯一且不可约；可约情形需用 packet 语言，多个特征层对应同一 $E$。原文 §5.4 p.73 只要求不可约，不要求酉。}
\intuit{不可约是「不能拆成更小的积木」；但不可约对象未必像素数那样有唯一分解。边界：不可约不等于酉、也不等于半单，几何侧没有正定性可用。}
\pitfall{把不可约当成酉表示：代数几何的 $\mathrm{Loc}$ 上没有紧实形式，酉性只是额外条件；也把不可约点当成普通点，忘记 $\mathrm{Loc}$ 是栈、点带自同构群。}
\end{word}

\begin{word}{分歧 / 非分歧}{Ramified / Unramified}{ramified-unramified}{§1 p.5}
\mathdef{设 $p\in C$，$C'=C\setminus\{p\}$。$C'$ 上的局部系统在 $p$ 附近给出单值共轭类；若可延拓到 $C$ 上（局部单值平凡）称非分歧，否则称在 $p$ 分歧。分歧又分 tame（一阶极点、正则奇点）与 wild（高阶极点，需 Stokes 数据）。数学上常改用带抛物结构的丛来记录单值。}
\mathex{$C=\mathbb P^1$ 非分歧时 $\pi_1$ 平凡，只有平凡系统；$E$ 带一个 tame 分歧点时 Hitchin 模空间是二维的，除三个点外纤维为光滑椭圆曲线（Frenkel §6.1 的例子）；$\mathbb P^1$ 上四个 tame 分歧点对应 Painlev\'e VI 的等单值形变。}
\physdef{物理侧：分歧点是插入 't Hooft 线算子的孔（puncture），非分歧对应无插入。tame 极点给出单极点式单值，wild 极点需 Stokes 矩阵；Nahm 极点给出 Dirichlet 型边界条件。}
\physex{近孔处 Hitchin 方程的解形如 $A=\alpha\,d\theta+\dots$、$\varphi=\beta\,dr/r-\gamma\,d\theta+\dots$；镜像对偶把参数按 $(\alpha,\eta)\mapsto({}^L\eta,-{}^L\alpha)$ 对调（Witten 2008）。}
\bridge{对应：孔 $\leftrightarrow$ Hecke 算子的插入点；分歧参数 $\leftrightarrow$ 扭化/边界参数；原文只处理非分歧情形，故本轨道各条锚点都默认 $C$ 紧致无孔。原文 §1 p.5。}
\intuit{非分歧是「曲线完整无缺」；分歧是「绕孔一圈单值不为 1」，像复变函数的支点，绕一圈换一张叶。边界：tame 不是「无分歧」，wild 更不能用单值完全刻画。}
\pitfall{把「非分歧」当成 $\pi_1$ 平凡：$C$ 本身可以不单连通；把 tame 与 wild 混同：wild 情形下局部系统由单值与 Stokes 数据共同决定，只看单值会漏参数。}
\end{word}

\begin{word}{主丛}{Principal Bundle}{principal-bundle}{§7.1 p.108}
\mathdef{$X$ 上的主 $G$-丛：$G$ 自由右作用、局部平凡（局部同构于 $U\times G$），转移函数 $g_{ij}:U_i\cap U_j\to G$，同构类由 $H^1(X,G)$ 分类；拓扑类由特征类 $c(E)\in H^2(X,\pi_1(G))$ 描述。向量丛就是 $G=\mathrm{GL}_n$ 的主丛配标准表示。}
\mathex{$G=GL_1$：主丛即线丛，$\mathrm{Pic}(C)=\mathbb Z\times\mathrm{Jac}(C)$；$G=SL_2$：行列式平凡给出次数为零，$\mathrm{Bun}_{SL_2}$ 连通；$C=\mathbb P^1$ 上单连通 $G$ 的主丛唯一（Bott 定理），故 $\mathrm{Bun}_G(\mathbb P^1)$ 退化为单点；一般 $G$ 的分量由 $\pi_1(G)$ 标号。}
\physdef{物理侧：规范场 $A$ 是主丛上的联络；开弦端点所带的 Chan-Paton 结构就是主丛（或其表示）；通量 $\xi(E)$ 是主丛的特征类，'t Hooft 线算子探测它。}
\physex{$G=U(1)$ 时主丛就是磁荷 $m$ 的线丛，Dirac 量子化给出 $m\in\mathbb Z$；$G=\mathrm{SO}(3)$ 的主丛由第二 Stiefel--Whitney 类区分，对应 't Hooft 通量。}
\bridge{对应：主丛 $=$ 规范场论的拓扑背景；特征类 $=$ 通量；GNO 对偶下 $G$ 的主丛通量与 $^LG$ 的表示配对。原文 §7.1 p.108。}
\intuit{主丛是「把 $G$ 粘贴到 $X$ 上的方式」：像给每点配一个 $G$，但没有典范的单位元可比较。边界：主丛不是向量丛，除非选一个表示把它变成向量丛。}
\pitfall{把主丛当成「纤维为 $G$ 的丛」的朴素图像：粘合时转移函数按共轭作用，两个不同的 $G$-丛可以有相同的底与纤维；也把特征类当成连续场，忘记它是拓扑不变量。}
\end{word}

\begin{word}{向量丛}{Vector Bundle}{vector-bundle}{§3.3 p.30}
\mathdef{$X$ 上向量丛：纤维为 $r$ 维向量空间、局部平凡、转移函数取值于 $\mathrm{GL}_r$，等价于局部自由 ${\cal O}_X$-模。复曲线 $C$ 上秩 $r$ 与次数 $\deg$ 给出斜率 $\mu=\deg/r$；稳定指任意真子丛满足 $\mu(F)<\mu(E)$，半稳定把不等号换成 $\le$。}
\mathex{$C=\mathbb P^1$ 上每个向量丛都同构于 $\oplus_i{\cal O}(a_i)$（Grothendieck 分解）；$E$ 为椭圆曲线时 $\deg 0$ 的线丛构成 $\mathrm{Pic}^0(E)\cong E$；$G=SL_2$ 的半稳定丛模空间维数为 $3g-3$。}
\physdef{物理侧：向量丛是膜上的 Chan-Paton 丛，开弦端点的指标住在其中；规范场的基本表示给出向量丛，因而电荷标号来自丛的秩与次数。}
\physex{磁特征膜带平坦 Chan-Paton 丛；Hecke 特征层的特征值 $V_E=E\times_{^LG}V$ 是 $C$ 上的平坦向量丛（Frenkel 的表述），其秩是 $V$ 的维数。}
\bridge{对应：主 $G$-丛加表示 $\rho$ 给出向量丛 $E\times_G V$；稳定向量丛给 BPS 态，Hecke 算子作用在向量丛上。原文 §3.3 p.30。}
\intuit{向量丛是「随位置转动的向量空间」：每点挂一个 $\mathbb C^r$，光滑过渡，整体可能扭转。边界：不是所有丛都平凡，整体扭转由特征类与次数度量。}
\pitfall{把稳定丛与「每个纤维稳定」混同：稳定性是整体条件，逐纤维无从谈起；把 $\deg$ 与 $r$ 当独立参数而忘记斜率 $\mu$，稳定性的比较只在同一斜率上进行。}
\end{word}

\begin{word}{模空间}{Moduli Space}{moduli-space}{§1 p.3}
\mathdef{模空间是对象按同构分类所得的空间；严格记录自同构要取模栈。本轨道涉及三个层次：稳定 $G$-丛模空间 $\mathcal M(G,C)$、所有 $G$-丛的栈 $\mathrm{Bun}_G(C)$、Hitchin 模空间 $\mathcal M_H(G,C)$（超凯勒）。$\mathrm{Bun}_G$ 的半稳定部分约化到 $\mathcal M$，而 $\mathcal M_H$ 含 $T^*\mathcal M$ 为 Zariski 开集。}
\mathex{$G=GL_1$：$\mathcal M=\mathrm{Pic}(C)$；$G=SL_2$：$\dim\mathcal M=3g-3$，$\dim\mathcal M_H=6g-6$；Hitchin 纤维化的底 $B=\oplus_i H^0(C,K^{d_i})$，$g\ge2$ 时 $\dim B=3g-3$。$\mathrm{Bun}_G$ 这类有奇点的模空间上，层的中间扩张由相交上同调（intersection cohomology）给出。}
\physdef{物理侧：模空间是真空模空间与靶空间；$\mathcal M_H(G)$ 是 $A$-模型靶空间，$\mathcal M_H(^LG)$ 是 $B$-模型靶空间，镜像对称联系两者，而 $\mathcal M_H$ 的复结构之间由 non-abelian Hodge 对应。}
\physex{$T^*\mathcal M$ 上的 $B$-膜给出 $\mathcal M$ 上的 ${\cal D}$-模，典范余各向同性膜给出这个翻译；$\mathcal M_H$ 是 $T^*\mathcal M$ 的部分紧化，故 ${\cal D}$-模是 $A$-膜的粗粒化。}
\bridge{对应：$\mathcal M(G,C)$ 上的扭化 ${\cal D}$-模 $\leftrightarrow$ $\mathcal M_H(G,C)$ 的 $A$-膜 $\leftrightarrow$ $\mathcal M_H(^LG,C)$ 的 $B$-膜；对偶群的 Hitchin 纤维化有同一底 $B$，这正是 SYZ 图像。原文 §1 p.3。}
\intuit{模空间是「把所有对象压成一个点」的空间，像按身高把一屋子人排成一条线。边界：非交换群要取栈，因为对象有自同构，点会变成 orbifold 点。}
\pitfall{把 $\mathcal M(G,C)$（稳定丛）当成 $\mathcal M_H(G,C)$（Hitchin 模空间）：前者不超凯勒，后者才装 $A/B$ 膜；把 $\mathrm{Bun}_G$ 当普通模空间：它是栈，Hecke 算子会把稳定丛跑出稳定部分。}
\end{word}

\begin{word}{Casimir 多项式}{Casimir Polynomial}{casimir-polynomial}{§4.3 p.53}
\mathdef{$\mathfrak g$ 上的 $G$-不变多项式环 $(\mathrm{Sym}\,\mathfrak g^*)^G=\mathbb C[c_1,\dots,c_r]$，次数 $d_1,\dots,d_r$（$r=\mathrm{rank}\,G$），如 $c=\mathrm{Tr}\,\varphi^k$；在曲线 $C$ 上 $\varphi\in H^0(C,K\otimes\mathrm{ad}E)$，故 $c_i(\varphi)\in H^0(C,K^{d_i})$，Hitchin 底 $B=\oplus_iH^0(C,K^{d_i})$。}
\mathex{$G=SL_2$：唯一基本 Casimir 是 $\mathrm{Tr}\,\varphi^2\in H^0(C,K^2)$，$\dim B=3g-3$；$G=SL_n$：$d_i=2,\dots,n$；$G=GL_1$：无 Casimir，$B$ 是单点。谱覆盖 $\det(\lambda-\varphi)=0$ 是 $B$ 上的纤维。}
\physdef{物理侧：Casimir 多项式是 Higgs 场的规范不变可观测量，$B$ 是质量/真空参数空间；在 Hitchin 可积系统中它们给出相互对易的哈密顿量。}
\physex{谱覆盖的系数就是 Casimir 多项式的值；$G$ 与 $^LG$ 的 Hitchin 底可自然等同（对偶群 Weyl 群相同、$d_i$ 相同），故两个 Hitchin 纤维化能逐纤维作 $T$-对偶。}
\bridge{对应：Casimir 多项式 $\leftrightarrow$ 谱覆盖系数 $\leftrightarrow$ 纤维上阿贝尔簇的参数；底的等同是镜像对称（SYZ）的关键。原文 §4.3 p.53。}
\intuit{Casimir 多项式像特征值的对称函数：模掉 Weyl 群后剩下的信息。边界：不变量不是特征值本身，Weyl 轨道没有典范排序，不能逐点还原 $\varphi$。}
\pitfall{把 Casimir 多项式当成 $\varphi$ 的特征值：后者依赖规范与排序，前者才是良定义函数；或以为 $d_i$ 与群无关：对偶群共享 $d_i$，但不同群不同（$GL_n$ 的 $d_i=2,\dots,n$）。}
\end{word}

\begin{word}{量子化}{Quantization}{quantization}{§1 p.4}
\mathdef{量子化：由经典相空间 $(M,\omega)$ 构造非交换代数或希尔伯特空间；形变量子化以 $\hbar$ 为参数把乘法变形为 Moyal 积，几何量子化取预量子化线丛并取极化；符号分次给出经典极限。}
\mathex{$T^*X$ 的形变量子化即 Weyl 代数，其层版本以 ${\cal D}_X$ 为原型；Hitchin 可积系统的量子化给出 $\mathrm{Bun}_G$ 上相互对易的微分算子（Beilinson--Drinfeld 的构造）；$G=GL_1$ 时量子化给出 $\mathrm{Pic}(C)$ 上的微分算子，与 Fourier--Mukai 变换相配。}
\physdef{物理侧：量子化是从经典场论到量子场论；原文 §12 的「量子几何 Langlands」研究把 $\hbar$ 引入几何 Langlands，而 Kapustin--Witten 对应把 $N=4$ 的 S-对偶翻译成几何 Langlands。}
\physex{四维 $N=4$ 的拓扑扭化带参数 $t$，S-对偶把 $t$ 变成 $\tau/|\tau|\,t$；量子化参数 $\hbar$ 与扭化参数相关，经典极限 $\hbar\to0$ 回到 Hitchin 可积系统。}
\bridge{对应：${\cal D}$-模是「余切丛量子化」在对象层面的体现，「$A$-膜 $\to$ ${\cal D}$-模」这一步就是量子化翻译；原文 §1 p.4 强调量子化与「取 ${\cal D}$-模」不是同义词，前者是经典--量子对应，后者是模层。}
\intuit{量子化像把可交换的函数代数换成不可交换的算子代数，如 $[x,p]=i\hbar$。边界：${\cal D}$-模不是量子希尔伯特空间，它没有内积、没有态，只有代数结构。}
\pitfall{把量子化当成「取 ${\cal D}$-模」的字面同义：量子化还含极化、内积与 $\hbar$ 的取值；也把 $\hbar$ 当纯几何参数而不追踪它在 S-对偶下的变换（$\hbar$ 与 $1/\hbar$ 的交换正是量子化的核心）。}
\end{word}

\begin{word}{非交换几何}{Noncommutative Geometry}{noncommutative-geometry}{§11 p.184}
\mathdef{非交换几何：用非交换代数 $A$ 代替函数环 ${\cal O}_X$ 来做几何；原型是 $T^*X$ 的形变量子化 $A_\hbar$，$[f,g]=\hbar\{f,g\}+\dots$，其层化与 ${\cal D}_X$ 同源（$\mathrm{gr}{\cal D}_X=\mathrm{Sym}\,T_X$）。}
\mathex{$\mathbb C^n$ 上 Weyl 代数 $A_n=\mathbb C\langle x_i,\partial_i\rangle$，$[\partial_i,x_j]=\delta_{ij}$；$\mathbb P^1$ 上的整体微分算子是 $U(\mathfrak{sl}_2)$ 的商；椭圆曲线上有非交换环面（$\theta$ 函数变形）。}
\physdef{物理侧：非交换几何来自背景 $B$-场：开弦世界叶在 $B$-场下坐标非交换，膜上的代数变成非交换代数；Kapustin 的论证说全纯辛流形非交换形变上的 $B$-膜范畴等价于其 $A$-膜范畴。}
\physex{常数 $B$-场给出 $[x,x]\sim\theta$ 的非交换环面；c.c. 膜的开弦代数 ${\cal D}_{K^{1/2}}$ 是交换情形的非交换替代，膜上的态就是它的模。}
\bridge{对应：${\cal D}$-模范畴 $=$ 余切丛非交换形变上的 $B$-膜范畴；这一步把「$A$-膜 $\leftrightarrow$ ${\cal D}$-模」的翻译升级为范畴等价（原文 §11，引 Kapustin 的论证）。原文 §11 p.184。}
\intuit{非交换几何是「点消失了，只剩函数代数」，像把空间换成它上面的算子代数。边界：这里的非交换来自余切丛的形变，不是把时空本身非交换化。}
\pitfall{把非交换几何当成「非交换时空」（$\theta$-形变时空）：原文指的是余切丛的非交换形变；也把 ${\cal D}_X$ 当成有限维非交换代数，它其实是层且有无穷维表示。}
\end{word}

\begin{word}{Fourier--Mukai 变换}{Fourier--Mukai Transform}{fourier-mukai}{§1 p.4}
\mathdef{设 $X,Y$ 光滑射影，核 $K\in D^b(X\times Y)$，定义函子 $\Phi_K:D^b(X)\to D^b(Y)$，$E\mapsto R\pi_{Y*}(\pi_X^*E\otimes^{\mathbb L}K)$；若 $\Phi_K$ 是范畴等价，称 Fourier--Mukai 等价。它是「用核做积分变换」的导出范畴版本。}
\mathex{椭圆曲线 $E$ 上 Poincaré 线丛给出 $D^b(E)\cong D^b(\hat E)$（$\hat E\cong E$），这是经典 FM 变换；阿贝尔簇同理；$G=GL_1$ 的几何 Langlands 就是 $\mathrm{Pic}(C)$ 上的 FM 变换（Laumon--Rothstein 证明）。}
\physdef{物理侧：FM 是 $T$-对偶的数学化——在环面纤维上把半径 $R$ 换成 $1/R$，交换动量与缠绕数，也就是交换电荷与磁荷；在 $B$-模型里 FM 对应沿纤维的 $T$-对偶。}
\physex{环面 $T^2$ 上 $T$-对偶交换 $(m,n)$ 荷；椭圆曲线的 FM 核就是 Poincaré 丛（通用线丛）；SYZ 镜像对称中纤维上的 $T$-对偶正是 FM。}
\bridge{对应：FM 把结构层/点送到向量丛（skyscraper 到平坦丛），正是「局部系统 $\leftrightarrow$ ${\cal D}$-模」的阿贝尔模型；Frenkel 称范畴几何 Langlands 为非交换 Fourier--Mukai。原文 §1 p.4。}
\intuit{FM 像「用核做积分变换」，傅里叶变换是它的原型：函数乘以核再积分。边界：必须用导出范畴与导出张量，否则核的 Tor 信息丢失。}
\pitfall{把 FM 当成普通前推/拉回函子：它由核定义，一般不保持结构层；忽略导出范畴与导出张量，会把非等价误当等价。}
\end{word}

\begin{word}{Kähler 商}{Kähler Quotient}{kahler-quotient}{§4.1 p.40}
\mathdef{辛流形 $(X,\omega)$ 带 Hamilton 群作用与 moment map $\mu:X\to\mathfrak g^*$，辛商为 $\mu^{-1}(0)/G$；若 $X$ 是 Kähler 且作用保复结构，商继承 Kähler 结构。超凯勒流形有三个复结构，须三个 moment map 同时取零再取商。}
\mathex{$\mathbb C^n/\!/U(1)=\mathbb P^{n-1}$；Hitchin 模空间由 ${\cal A}\times\Omega^{1,0}(C,\mathrm{ad}E)$ 作超凯勒商得到；$T^*X$ 的辛商给出余切丛的约化，是 ${\cal D}$-模与非交换几何的几何背景。}
\physdef{物理侧：Kähler 商 $=$ 先解 D-项方程（moment map 取零）再模规范群；Hitchin 方程 $F-\varphi\wedge\varphi=0$、$D\varphi=D\star\varphi=0$ 正是超凯勒商的三组方程，解空间模规范群得 $\mathcal M_H$。}
\physex{$\mathcal M_H$ 是超凯勒商，故自然带 $I,J,K$ 三个复结构；$A$-模型用辛形式 $\omega_K$，$B$-模型用复结构，两者由镜像对称联系。}
\bridge{对应：D-项方程 $\leftrightarrow$ moment map；F-项方程（全纯性）$\leftrightarrow$ 复结构；四维 $N=4$ 的真空模空间 $=$ 超凯勒商 $=$ $\mathcal M_H$。原文 §4.1 p.40。}
\intuit{Kähler 商是「先挑出零动量面，再按对称群压扁」，像把旋转对称的球面做成一片。边界：超凯勒商要三个 moment map 都取零，只取一个会丢掉两个复结构。}
\pitfall{把 Kähler 商与 GIT 商混同：两者只在特定稳定性条件下一致；把超凯勒商当普通辛商：后者丢掉 $I,K$ 复结构与超凯勒度量。}
\end{word}

\begin{word}{通量}{Flux}{flux}{§7.1 p.103}
\mathdef{主 $G$-丛 $E$ 的通量是特征类 $\xi(E)\in H^2(C,\pi_1(G))$；离散电、磁通量分别为 $e\in X^*(T)$（weight）与 $m\in X_*(T)$（coweight），Dirac 量子化要求配对 $\langle e,m\rangle\in\mathbb Z$。$G=U(1)$ 时通量就是第一 Chern 类 $\in\mathbb Z$。}
\mathex{$C=\mathbb P^1$ 上单连通 $G$ 的主丛唯一，故无通量；$E$ 为椭圆曲线时 $\mathrm{Bun}_G$ 的分量由 $\pi_1(G)$（或次数）标号；$G=U(1)$ 时通量 $m\in\mathbb Z$，$G=\mathrm{SO}(3)$ 时第二 Stiefel--Whitney 类给出 $\mathbb Z_2$ 通量。}
\physdef{物理侧：'t Hooft 线算子产生磁通量 $m$，Wilson 线对应电通量 $e$；$N=4$ 中 $(e,m)$ 在 $SL(2,\mathbb Z)$ 下成轨道，S-对偶交换 $e$ 与 $m$；GNO 对偶交换两个格。}
\physex{'t Hooft 算子按 $^LG$ 的表示标号，Wilson 圈按 $G$ 的表示标号；磁单极的 Dirac 量子化 $m\cdot e\in\mathbb Z$ 就是两个格的对偶配对。}
\bridge{对应：电通量格 $X_*(T)$ 与磁通量格 $X^*(T)$ 在 GNO 下互换，正是 $^LG$ 交换 weight/coweight 格；通量 $\xi(E)$ 决定 $\mathrm{Bun}_G$ 的连通分支。原文 §7.1 p.103。}
\intuit{通量是「绕不过去的洞」：像水管里绕一圈的水流，局部场强可以为零而整体仍有流量。边界：通量是拓扑类，不是场强密度。}
\pitfall{把 flux 当成电磁场的场强密度 $E,B$（原文 §7.1 明确不是）；把电、磁通量当成同一格里的数：它们是互相对偶的两个格，配对才给出整数。}
\end{word}

\begin{word}{de Rham 微分}{de Rham Differential}{de-rham-differential}{§4.1 p.45}
\mathdef{流形 $X$ 上的外微分 $d:\Omega^k\to\Omega^{k+1}$，$d^2=0$；de Rham 上同调 $H^k_{dR}=\ker d/\mathrm{im}\,d$。复流形上 $d=\partial+\bar\partial$；原文用 $\delta$ 记这个微分（§4.1 p.45），它是作用在形式上的一个实算子。}
\mathex{$C$ 上 $H^1(C,\mathbb R)$ 的维数是 $2g$，由闭形式模恰当形式给出；${\cal D}_X$ 的一阶符号层正是 $T_X$，即 $d$ 的符号；$H^1_{dR}(C)=H^{1,0}\oplus H^{0,1}$，即 $d=\partial+\bar\partial$ 的分解；$X=\mathbb P^1$ 时 $H^0_{dR}=\mathbb R$、$H^2_{dR}=\mathbb R$。}
\physdef{物理侧：$d$ 是规范变换与 BRST 算子的一部分；场强 $F=dA$ 是 $d$ 作用在联络上；平坦联络即 $d_A^2=0$。拓扑扭化后 BRST 荷的经典部分常是 $d$ 或 $\bar\partial$。}
\physex{$B$-模型的 BRST 荷以 $\bar\partial$ 为主，$A$-模型以 $d$ 为主；扭化理论中 $Q$ 是超荷的标量部分，其经典近似给出 de Rham 微分。}
\bridge{对应：${\cal D}$-模的联络 $\nabla$ 是 $d$ 的「丛版本」（$\nabla:M\to M\otimes\Omega^1$ 在截面上就是 $d$）；de Rham 栈上的层即 ${\cal D}$-模。原文 §4.1 p.45 的 $\delta$ 就是 $d$。}
\intuit{$d$ 是「把函数变微分」的算子，像微积分里的全微分。边界：de Rham 上同调等于奇异上同调要经 de Rham 定理，不能直接当成拓扑定义。}
\pitfall{把 $d$ 与 Dolbeault 算子 $\bar\partial$ 混同：复结构下 $d=\partial+\bar\partial$，$B$-模型只用 $\bar\partial$；把 $d$ 与协变导数 $d_A$ 混同：后者带联络，曲率非零时不平方为零。}
\end{word}

\begin{word}{消灭定理}{Vanishing Theorem}{vanishing-theorem}{§5.1 p.65}
\mathdef{消灭定理：某些层或上同调为零。例：Serre 消灭（扭转到足够正时高次上同调为零）、Kodaira 消灭（$K_X\otimes L$ 的 $H^{>0}$ 为零，$L$ 丰沛）、Borel--Weil--Bott（旗流形上线丛的上同调集中于一次）。}
\mathex{$X=\mathbb P^1$：$H^1(\mathbb P^1,{\cal O}(n))=0$ 对 $n\ge-1$，而 $H^0({\cal O}(-1))=0$；仿射簇上相干层的高次上同调为零；$G/B$ 上 Borel--Weil 把表示实现为线丛的上同调。}
\physdef{物理侧：$A$-模型态空间是常值映射的上同调，用消灭定理可算；$B$-模型态空间是 Ext 群或全纯形式，消灭定理让开弦谱退化。原文 §5.1 p.65 用它说明 $A$-模型态空间。}
\physex{对 $\mathbb P^1$ 这类靶空间，非平凡世界叶的贡献被度数条件排除，态空间只剩常值映射的上同调；环面上则按同调类分层。}
\bridge{对应：$H^{>0}(C,\mathrm{ad}E)=0$ 这类条件 $\leftrightarrow$ 无穷小形变无阻、真空孤立；消灭定理在两侧的角色都是「把态空间简化成有限维上同调」。原文 §5.1 p.65。}
\intuit{消灭定理像「大部分高阶项自动消失」，像分部积分把边界项杀掉。边界：需要正性或度数条件，条件不满足时高次上同调非零。}
\pitfall{把消灭定理当成无条件成立：正性/度数条件是前提；把上同调消失误读为「层本身为零」：层可以非零而上同调为零（如 ${\cal O}(-1)$ 在 $\mathbb P^1$ 上）。}
\end{word}
